\section{The original transitions between ideal-power cohomologies}

\begin{lemma}[Literal original ideal-power transition maps]
\label{lem:fc-rees-cech-degree-transitions}
\leanok
\lean{EnvelopeProperty.FcReesCechDegreeTransitions.termLower, EnvelopeProperty.FcReesCechDegreeTransitions.termLower_coface, EnvelopeProperty.FcReesCechDegreeTransitions.termLower_differential, EnvelopeProperty.FcReesCechDegreeTransitions.degreeLower, EnvelopeProperty.FcReesCechDegreeTransitions.degreeLower_self, EnvelopeProperty.FcReesCechDegreeTransitions.degreeLower_comp, EnvelopeProperty.FcReesCechDegreeTransitions.degreeMultiply_lower, EnvelopeProperty.FcReesCechDegreeTransitions.homologyLower, EnvelopeProperty.FcReesCechDegreeTransitions.homologyLower_self, EnvelopeProperty.FcReesCechDegreeTransitions.homologyLower_comp, EnvelopeProperty.FcReesCechDegreeTransitions.homologyMultiply_lower}
\uses{lem:fc-rees-cech-degree-complex,lem:fc-rees-cech-homogeneous-multiplication}
For $n\leq m$, the original inclusions of ideal powers give cochain maps
from the literal degree-$m$ complex to the literal degree-$n$ complex.
They retain the original underlying section on every affine intersection,
are the identity for $m=n$, and compose according to the original order.
If $a\in I^e$, homogeneous multiplication followed by lowering to degree
$n$ equals the original coefficient $a$ times the original lowering map
from degree $m$ to degree $n$.  These identities hold on every
cohomology group as well.
\end{lemma}
\begin{proof}
The inclusion $I^m\subseteq I^n$ retains each original coefficient.
Applying it on every affine intersection therefore commutes with every
restriction and every ordered coface.  Taking the alternating sums gives
a cochain map.  The same coefficient identity proves its identity and
composition laws.  Multiplication by $a$ sends a section $s$ to
$(a|_V)s$ on an intersection $V$.  Forgetting its higher ideal-power
membership leaves exactly the original scalar multiple of the
underlying section.  This proves the asserted equality of entire
cochain maps.  The homology functor preserves compositions, identities,
and the original base scalar multiplication, so it preserves all these
equalities.
\end{proof}
